\chapter{Evaluation}
\label{chap:evaluation}

\section{Training loss}
\label{sec:training-loss}

A network that spreads its probability evenly over all 128 bins has a loss
of $\ln 128 \approx 4.85$, which is a reference point for the values below.
The two training runs reached almost the same best validation loss: 2.4857
for v1 and 2.4855 for v2. Both
losses fell from about 4.0 after the first 1,000 steps, and in both runs the
validation loss still varied by around 0.02 from one evaluation to the next
in the final 10,000 steps. A difference of 0.0002 is far inside that
variation, so by validation loss the two models are equal.
In games, v2 plays stronger than v1. Validation loss is
therefore a weak guide to playing strength, and Instinct's strength is
measured in games instead.

Figure~\ref{fig:loss} shows the loss of both runs. Both have short spikes.
The validation loss of v1 jumps from about 2.85 to 3.96 at step 64,000, and
that of v2 rises to 3.44, 3.32 and 3.11 at steps 71,000, 187,000 and
219,000. At each spike the training loss jumps as well, so the model itself
got worse for a moment, and by the next evaluation, 1,000 steps later, the
loss is back on its curve. The cause was not investigated. Instinct is
trained without learning-rate warm-up, without learning-rate decay and
without gradient clipping; whether any of these would have prevented the
spikes was not tested. Keeping the checkpoint with the lowest validation
loss ensures that a run cannot end on a spike.

\begin{figure}[htbp]
  \centering
  \begin{tikzpicture}
  \begin{axis}[
      width=\textwidth, height=0.5\textwidth,
      xlabel={Training step}, ylabel={Loss},
      xmin=0, xmax=270000, ymin=2.4, ymax=4.1,
      scaled x ticks=false,
      xtick={0,50000,100000,150000,200000,250000},
      xticklabels={0,50k,100k,150k,200k,250k},
      legend pos=north east, legend cell align=left,
      grid=major, grid style={gray!20},
    ]
    \addplot[blue!35, thin] table[x=step, y=train] {figures/data/loss-v1.dat};
    \addlegendentry{v1 training}
    \addplot[blue, thick] table[x=step, y=val] {figures/data/loss-v1.dat};
    \addlegendentry{v1 validation}
    \addplot[orange!45, thin] table[x=step, y=train] {figures/data/loss-v2.dat};
    \addlegendentry{v2 training}
    \addplot[orange!90!black, thick] table[x=step, y=val] {figures/data/loss-v2.dat};
    \addlegendentry{v2 validation}
  \end{axis}
\end{tikzpicture}

  \caption{Training and validation loss of v1 and v2 over the single pass. The
    training loss is that of the batch at each logged step; the validation loss
    is computed on 200,000 samples every 1,000 steps.}
  \label{fig:loss}
\end{figure}

\section{Vision}
\label{sec:vision}

The classifier named every validation crop correctly from the first epoch on.
